% Complete distinct Stirling route and Durrett Section 3.1 exercises.
\section{Stirling's formula and the de Moivre--Laplace theorem}
\label{sec:annex-p15-de-moivre-laplace-and-stirling}

The proofs follow Durrett's factorial route
\cite[Section~3.1, pp.~113--115]{DurrettPTEAuthor2019},
with the full Feller Stirling referral and its constant argument
\cite[II.9 and VII.2, pp.~51--54, 179--181]{Feller1968ProbabilityI}.

\begin{annexitem}{annex:p15.stirling-log-tools}{thm:p15-stirling-log-tools}{Finite remainders and both source product statements}
\begin{proof}
For \(|t|\leq1/2\),
\[
 \log(1+t)-t=\int_0^t\frac{-v}{1+v}\,dv,
\]
whose absolute value is at most
\(\int_0^{|t|}2v\,dv=t^2\).
For \(0<y<1\), integrate the finite geometric identity
\[
 \frac1{1-v^2}=\sum_{j=0}^m v^{2j}
                   +\frac{v^{2m+2}}{1-v^2}
 \quad(0\leq v\leq y).
\]
Its integrated remainder is nonnegative and at most
\(y^{2m+3}/((2m+3)(1-y^2))\), which tends to zero.
Since
\(\int_0^y(1-v^2)^{-1}dv=\tfrac12\log((1+y)/(1-y))\),
division by \(y>0\) proves the asserted series. Thus the
logarithmic expansion used in the source is established
with its limiting interchange justified.

For the scalar limit, \(a_n>0\) eventually, and the first
estimate gives
\[
 |a_n\log(1+c_n)-a_nc_n|
                 \leq |a_nc_n|\,|c_n|\longrightarrow0.
\]
The bases are eventually positive, so exponentiation proves
the result for real exponents, including \(c_n=0\).
For the triangular version, put
\(M_n=\max_{j\leq n}|c_{n,j}|\).
Eventually \(M_n\leq1/2\), and
\[
 \left|\sum_{j=1}^n\log(1+c_{n,j})
                  -\sum_{j=1}^nc_{n,j}\right|
 \leq\sum_{j=1}^n|c_{n,j}|^2
 \leq M_n\sum_{j=1}^n|c_{n,j}|\longrightarrow0.
\]
Exponentiation gives the source exercise's full conclusion.
\end{proof}
\end{annexitem}

\begin{annexitem}{annex:p15.stirling-tail-premises}{thm:p15-stirling-tail-premises}{The complete two source tail estimates}
\begin{proof}
There are \(\binom Nk\) choices of the \(k\) successful
trials, and each prescribed success/failure pattern has
probability \(p^kq^{N-k}\) by independence.
The resulting disjoint events partition the sample space,
which proves both the mass formula and its normalization.
Writing \(B_N\) as a sum of the success indicators gives
mean \(Np\); the individual variance is \(p-p^2=pq\),
and independent variance addition gives \(Npq\).

Write \(b(k)=b(k;N,p)\). For \(0\leq k<N\),
\[
 \frac{b(k+1)}{b(k)}=\frac{(N-k)p}{(k+1)q}.
\]
If the integer \(r>Np\), then for \(k\geq r\) this is at
most \((N-r)p/(rq)=1-(r-Np)/(rq)<1\).
If \(r=N\), the tail has its single term and the same
bound has ratio zero.
Comparison with the geometric series therefore gives
\[
 \mathbb P(B_N\geq r)\leq b(r)\frac{rq}{r-Np}.
\]
The masses are nonincreasing for integers
\(k\geq\lceil Np\rceil\): the displayed ratio is below
one because \(k+1\geq Np+1>(N+1)p\).
Every term from \(\lceil Np\rceil\) through \(r\) is thus
at least \(b(r)\). There are
\(r-\lceil Np\rceil+1>r-Np\) such terms.
Their sum is at most one, so \(b(r)\leq1/(r-Np)\).
This gives Feller's second bound.
The variable \(N-B_N\) counts failures and has the
binomial law with probabilities \(q,p\) exchanged.
Applying the upper-tail bound to its threshold \(N-r>Nq\)
gives the stated lower-tail bound. Every division used a
strictly positive denominator.

The evenness and previously proved normalization of \(\phi\)
give \(\Phi(0)=1/2\). For the Gaussian bounds, differentiation gives
\(\phi'(x)=-x\phi(x)\). For \(x>0\),
\[
 -\frac{d}{dx}\left(\frac{\phi(x)}x\right)
          =(1+x^{-2})\phi(x)>\phi(x),
\]
and
\[
 -\frac{d}{dx}\left((x^{-1}-x^{-3})\phi(x)\right)
          =(1-3x^{-4})\phi(x)<\phi(x).
\]
Both differentiated primitives tend to zero at positive
infinity, by exponential decay. Integrating from \(x\)
to infinity proves the strict upper and lower bounds.
The lower bound is allowed to be negative for small \(x\).
For \(x\to\infty\), division by \(\phi(x)/x>0\)
squeezes the tail ratio between \(1-x^{-2}\) and one,
proving the asymptotic assertion.
\end{proof}
\end{annexitem}

\begin{annexitem}{annex:p15.stirling-with-remainder}{thm:p15-stirling-with-remainder}{The full Feller--Robbins argument and both constant determinations}
\begin{proof}
We first retain the source's initial comparison.
Monotonicity of \(\log x\) gives
\[
 \int_0^n\log x\,dx<\log n!<
                    \int_1^{n+1}\log x\,dx,
\]
or \(n\log n-n<\log n!<(n+1)\log(n+1)-n\).
The improper integral at zero is finite because
\(x\log x\to0\) as \(x\downarrow0\).
Define, as in Feller,
\[
 d_n=\log n!-(n+\tfrac12)\log n+n,\qquad
 \Delta_n=d_n-d_{n+1}.
\]
Cancellation gives
\[
 \Delta_n=(n+\tfrac12)\log(1+1/n)-1.
\]
Put \(y=(2n+1)^{-1}\), so
\((n+1)/n=(1+y)/(1-y)\).
The proved logarithmic series yields the exact expansion
\[
 \Delta_n=\sum_{j=1}^\infty
                 \frac{1}{(2j+1)(2n+1)^{2j}}.
\]
Every term is positive. Bounding each denominator \(2j+1\)
by three, with strict inequality beyond the first term, gives
\[
 0<\Delta_n<
 \frac1{3((2n+1)^2-1)}
     =\frac1{12n}-\frac1{12(n+1)}.
\]
Hence \(d_n\) decreases, while its total possible decrease
from any starting index is bounded by the telescoping series.
It has a finite limit \(C\).
For the lower companion bound, retain the first positive
term and observe, for every integer \(n\geq1\),
\[
 \Delta_n>\frac1{3(2n+1)^2}>
           \frac1{12n+1}-\frac1{12(n+1)+1}.
\]
The second strict inequality follows after multiplication
by positive denominators from
\((12n+1)(12n+13)-36(2n+1)^2=24n-23>0\).
The same comparisons show that \(d_n-(12n)^{-1}\) increases
and \(d_n-(12n+1)^{-1}\) decreases, as in the source argument.
Summing the two comparisons from \(n\) to infinity shows that
\[
 r_n=d_n-C=\sum_{j=n}^\infty\Delta_j,\qquad
 \frac1{12n+1}<r_n<\frac1{12n}.
\]
In particular \(r_n\downarrow0\) strictly, and with \(A=e^C>0\)
we have the exact, as yet unnormalized formula
\[
 n!=A\,n^{n+1/2}e^{-n}e^{r_n}.
\]

We now discharge the constant referral by Feller's
probability-mass argument. This argument does not use a
previously established normal limit.
For the fair-sign walk define
\[
 p_n(k)=4^{-n}\binom{2n}{n+k},\qquad -n\leq k\leq n.
\]
The formula with unknown \(A\) implies
\[
 p_n(0)=\frac{\sqrt2}{A\sqrt n}\,e^{r_{2n}-2r_n}
                          \sim\frac{\sqrt2}{A\sqrt n}.
\]
The exact factorial ratio is
\[
 \frac{p_n(k)}{p_n(0)}
   =\prod_{j=1}^k\frac{n-j+1}{n+j}\quad(0\leq k\leq n).
\]
Fix \(0<M<\infty\). For \(k\leq M\sqrt n\), the logarithmic
error bound gives
\[
 \log\frac{p_n(k)}{p_n(0)}
    =-\frac{k^2}{n}+E_{n,k},\qquad
 |E_{n,k}|\leq\frac2{n^2}\sum_{j=1}^k j^2
                         \longrightarrow0
\]
uniformly in that range. Symmetry handles negative \(k\).
This is the central-ratio argument in the cited VII.2
passage with the constant still undetermined.
With \(h_n=2/\sqrt{2n}\) it gives, uniformly on every
bounded interval of \(x=k h_n\),
\[
 \frac{p_n(k)}{h_n}\longrightarrow A^{-1}e^{-x^2/2}
                  =c\phi(x),\qquad c=\frac{\sqrt{2\pi}}A.
\]
Here the uniform limit notation means that the supremum
of the difference over the indicated lattice points tends
to zero.
For fixed \(R>0\), summing the uniform estimate gives
\[
 \mathbb P(0<S_{2n}/\sqrt{2n}\leq R)
       \longrightarrow c\int_0^R\phi(x)\,dx.
\]
For completeness, the corresponding Riemann sum converges
without any limit theorem: for decreasing
\(g(x)=e^{-x^2/2}\) on \([0,R]\), its upper and lower sums
on the mesh \(h_n\) differ by at most \(h_n g(0)\);
the last partial interval has length less than \(h_n\).
Multiplication by the fixed \(A^{-1}\) preserves the limit.

Since \(p_n(0)\to0\), sign symmetry gives
\(\mathbb P(S_{2n}>0)=(1-p_n(0))/2\to1/2\).
Feller's binomial tail premise also supplies the required
uniformly vanishing outer tail as \(R\to\infty\).
Indeed \(B_{2n}=(S_{2n}+2n)/2\) has \(p=q=1/2\).
For the event \(S_{2n}/\sqrt{2n}>R\), use
\(r=\lfloor n+R\sqrt{n/2}\rfloor+1\).
For sufficiently large \(n\), \(r\leq2n\) and
\(r-n>R\sqrt{n/2}\), so the source bound gives a tail
probability at most \(2/R^2\), and hence at most the
source's more conservative \(4/R^2\).
Consequently
\[
 \tfrac12-4/R^2\leq c(\Phi(R)-\tfrac12)\leq\tfrac12.
\]
The other source premise proved above gives
\[
 c\left(\tfrac12-\frac{\phi(R)}R\right)
       <c(\Phi(R)-\tfrac12)<c/2.
\]
Letting \(R\to\infty\) forces \(c/2=1/2\), so \(A=\sqrt{2\pi}\).
This completes Feller's deferred constant argument by its
own central-ratio, interval-mass and tail comparisons,
without presuming the normalized de Moivre--Laplace conclusion.

There is also the independent determination mentioned in
the body. The cosine-integral proof in
Theorem~\ref{thm:p15-weak-elementary-asymptotics}
established \(\sqrt n\,4^{-n}\binom{2n}{n}\to1/\sqrt\pi\)
using neither factorial asymptotics nor any normal-limit
theorem. The unknown-\(A\) identity above gives the same
limit as \(\sqrt2/A\); equating them again yields
\(A=\sqrt{2\pi}\). This second route is logically independent
of the probability-mass constant determination.

Finally we prove the refinement with an explicit bound.
Put \(s_n=1/(12n)-1/(360n^3)\) and \(x=1/n\in(0,1]\).
Integrating a finite geometric series for \(1/(1+x)\)
gives the alternating logarithmic expansion through degree
six with a remainder between zero and \(x^7/7\).
Multiplication by \(1/x+1/2\) shows
\[
 \Delta_n=
 \frac{x^2}{12}-\frac{x^3}{12}
       +\frac{3x^4}{40}-\frac{x^5}{15}+e_n,\qquad
 |e_n|\leq\frac{x^6}{3}.
\]
The bound follows directly from the unretained term
\(-x^6/12\) and the remainder contribution, whose absolute
sum is at most \(x^6(1/12+3/14)<x^6/3\).
On the other hand the exact rational expression is
\[
 s_n-s_{n+1}
  =\frac{x^2}{12(1+x)}
       -\frac{x^4}{360}\frac{3+3x+x^2}{(1+x)^3}.
\]
The first term has expansion
\((x^2-x^3+x^4-x^5)/12\) with remainder of absolute
value at most \(x^6/12\). The identity
\[
 \frac{3+3x+x^2}{(1+x)^3}-(3-6x)
             =\frac{10x^2+15x^3+6x^4}{(1+x)^3}
\]
bounds the second remainder by \(31x^6/360\).
Thus \(s_n-s_{n+1}\) has the same four displayed terms
as \(\Delta_n\), with remainder at most \(61x^6/360\).
It follows that
\[
 |\Delta_n-(s_n-s_{n+1})|\leq n^{-6}.
\]
Sum from \(n\) to infinity, using \(s_j\to0\), to obtain
\[
 |r_n-s_n|\leq\sum_{j=n}^\infty j^{-6}
       \leq n^{-6}+\int_n^\infty t^{-6}dt
       \leq\frac6{5n^5}.
\]
This proves the claimed refinement as well as the full
Stirling formula and both Robbins inequalities.
\end{proof}
\end{annexitem}

\begin{annexitem}{annex:p15.stirling-local-rademacher}{thm:p15-stirling-local-rademacher}{The source factorial factors and a uniform local error}
\begin{proof}
Counting signs gives
\(\mathbb P(S_{2n}=2k)=4^{-n}(2n)!/((n+k)!(n-k)!)\).
For \(-n<k<n\), all factorial indices are positive, so
substitution of the exact Stirling formulas yields
\[
 \mathbb P(S_{2n}=2k)
  =(\pi n)^{-1/2}
    (1+u)^{-n-k}(1-u)^{-n+k}
    (1+u)^{-1/2}(1-u)^{-1/2}e^{R_{n,k}},
 \qquad u=k/n.
\]
This is Durrett's factorization with its remainder retained.
Taking logarithms of the middle powers proves the asserted
exact formula involving \(H(u)\).
The source also separates the two leading powers into
\[
 (1-u^2)^{-n}(1+u)^{-k}(1-u)^k.
\]
If \(x_{n,k}=2k/\sqrt{2n}\to x\), then \(u\to0\) and
\(nu^2=x_{n,k}^2/2\). The logarithms of these three
factors tend respectively to \(x^2/2,-x^2/2,-x^2/2\).
For the last two claims, the logarithmic remainder is
bounded by \(n|u|^3\to0\); for the first use
\(|\log(1-u^2)+u^2|\leq u^4\) eventually.
This proves the source product limits also when \(k\)
is negative or zero, without an unjustified positive-exponent
assumption. The square-root factors tend to one, and
\(R_{n,k}\to0\), giving the stated source limit.

To establish the uniform bound needed for summation,
differentiate \(H\):
\[
 H(0)=H'(0)=0,\qquad H''(u)=\frac2{1-u^2}.
\]
It is even, and for \(|u|\leq1/2\),
\[
 0\leq H(u)-u^2
 =\int_0^{|u|}(|u|-t)\frac{2t^2}{1-t^2}\,dt
 \leq\frac{2u^4}{9}.
\]
Also
\[
 0\leq-\tfrac12\log(1-u^2)\leq\tfrac23u^2,\qquad
 |R_{n,k}|\leq\frac1{24n}
      +\frac1{12(n+k)}+\frac1{12(n-k)}
 \leq\frac3{8n}.
\]
The first bound follows by integrating \((1-t)^{-1}\)
up to \(u^2\leq1/4\); the last follows from
\(n\pm k\geq n/2\).
With \(u=x_{n,k}/\sqrt{2n}\), these inequalities give
exactly
\[
 \left|-n(H(u)-u^2)-\tfrac12\log(1-u^2)+R_{n,k}\right|
 \leq \frac{x_{n,k}^4}{18n}
        +\frac{x_{n,k}^2}{3n}+\frac3{8n}.
\]
The condition \(n\geq2A^2\) ensures \(|u|\leq1/2\)
throughout \(|x_{n,k}|\leq A\), proving the displayed
uniform logarithmic remainder and hence the even local
limit. In particular the relative error there is at most
\(e^{K_A/n}-1\).

For odd indices, independence of the last sign gives
\[
 \mathbb P(S_{2n+1}=2k+1)
     =\tfrac12\{\mathbb P(S_{2n}=2k)
                         +\mathbb P(S_{2n}=2k+2)\}.
\]
If \(|2k+1|/\sqrt{2n+1}\leq A\), the two corresponding
even scaled arguments \(2k/\sqrt{2n}\) and
\((2k+2)/\sqrt{2n}\) differ from
\((2k+1)/\sqrt{2n+1}\) by a quantity tending uniformly
to zero. Explicitly its absolute value is at most
\[
 A\left(\sqrt{\frac{2n+1}{2n}}-1\right)+\frac1{\sqrt{2n}}.
\]
They lie in a fixed enlarged bounded interval eventually.
Use the even uniform estimate there, uniform continuity
of \(\phi\), and \(\sqrt{(2n+1)/(2n)}\to1\) in this
two-term formula. This proves the stated odd local limit.
Wrong-parity points have no sign pattern and probability
zero; the extreme points each have just one pattern,
of probability \(2^{-N}\).
\end{proof}
\end{annexitem}

\begin{annexitem}{annex:p15.stirling-de-moivre-laplace}{thm:p15-stirling-de-moivre-laplace}{Rounding, Riemann sums, parity and endpoints}
\begin{proof}
If exactly \(j\) of the \(N\) signs are positive, then
\(S_N=2j-N\). Solving
\(a\sqrt N\leq2j-N\leq b\sqrt N\), imposing integer \(j\)
and intersecting with \(0\leq j\leq N\), gives exactly
the displayed lower and upper indices. This proves the
finite formula, including empty intervals.

The scaled support has spacing \(h_N=2/\sqrt N\).
For fixed finite \(a<b\), the uniform local theorem gives
\[
 \mathbb P(a\leq S_N/\sqrt N\leq b)
   =\sum_{x\in[a,b]\cap(\mathcal L_N/\sqrt N)}
                          h_N\{\phi(x)+e_N(x)\},
\]
where the supremum of \(|e_N(x)|\) on the interval tends
to zero. There are at most \((b-a)/h_N+2\) terms, so the
total error is at most
\((b-a+2h_N)\sup|e_N|\to0\).
The remaining sum tends to \(\int_a^b\phi\):
uniform continuity on an enlarged compact interval bounds
the error on each full lattice cell by its length times
the modulus of continuity of \(\phi\) at \(h_N\);
the two truncated end cells have total length at most
\(2h_N\) and bounded integrand.
This proves the source summation step with uniform error.

It also proves the result by the source's explicitly
separated parity route. First use the even version alone
in the preceding Riemann sums. Independent signs give
\(\mathbb E S_N^2=N\), so
\(\mathbb P(|S_N|/\sqrt N>A)\leq A^{-2}\).
Cutting a CDF at a finite left endpoint \(-A\), applying
the bounded-interval even limit, and letting \(A\to\infty\)
proves the even CDF limit.
For \(d_n=\sqrt{2n/(2n+1)}\), the exact relation
\(S_{2n+1}=S_{2n}\pm1\) yields
\[
 \mathbb P\!\left(\frac{S_{2n}}{\sqrt{2n}}
                \leq x/d_n-1/\sqrt{2n}\right)
 \leq\mathbb P(S_{2n+1}/\sqrt{2n+1}\leq x)
 \leq
 \mathbb P\!\left(\frac{S_{2n}}{\sqrt{2n}}
                \leq x/d_n+1/\sqrt{2n}\right).
\]
Both thresholds tend to \(x\). Bounding them between
\(x-\eta,x+\eta\), applying the even CDF limit and then
letting \(\eta\downarrow0\), gives \(\Phi(x)\) by its
continuity. Thus both the all-parity local route and the
source's even-limit/odd-step route are complete.

For any fixed finite endpoint \(x\), its point probability
is either zero or, by the local estimate in a bounded
window containing \(x\), at most a constant times
\(N^{-1/2}\). It tends to zero. Changing strict to closed
endpoints therefore affects the interval probability by
at most two vanishing terms; the singleton case follows
as well. The CDF limit holds for every real \(x\).
It and the already proved tail estimate give the infinite
endpoint interval versions, or these follow by finite
cutoffs and then letting the cutoffs escape to infinity.
At \(a=-\infty,b=+\infty\), both sides are one.
\end{proof}
\end{annexitem}

\begin{annexitem}{annex:p15.stirling-poisson-premises}{thm:p15-stirling-poisson-premises}{Exact normalization and independent sums}
\begin{proof}
For \(\lambda>0\), the exponential series
\(\sum_{k\geq0}\lambda^k/k!=e^\lambda\) gives normalization.
This identity itself follows from Taylor's formula for
\(e^\lambda\): its remainder after degree \(m\) is bounded
by \(e^\lambda\lambda^{m+1}/(m+1)!\).
The last factor tends to zero, because its consecutive
ratio is \(\lambda/(m+2)\), eventually at most \(1/2\).
Thus no exchange of an infinite derivative and a sum is
needed. Shifting indices in convergent nonnegative series gives
\[
 \mathbb EX=\lambda,\qquad
 \mathbb E[X(X-1)]=\lambda^2,\qquad
 \mathbb EX^2=\lambda^2+\lambda.
\]
Hence the variance is \(\lambda\).
At \(\lambda=0\), all assertions follow from the point
mass convention.

For \(k\geq0\), independence and the fact that summands
are nonnegative integers give the finite sum
\[
 \mathbb P(T_n=k)
     =e^{-n}\sum_{k_1+\cdots+k_n=k}\frac1{k_1!\cdots k_n!}.
\]
There are \(n^k\) assignments of \(k\) labelled objects
to \(n\) labelled boxes. Those with prescribed box counts
\(k_1,\ldots,k_n\) number \(k!/(k_1!\cdots k_n!)\),
as seen by choosing the contents of each box successively.
Summing this partition of the assignments gives the
finite multinomial identity and hence
\(\mathbb P(T_n=k)=e^{-n}n^k/k!\).
For \(k=0\) the identity is the empty assignment and is
still valid. The law is therefore Poisson of parameter
\(n\), giving the displayed mean and variance.
\end{proof}
\end{annexitem}

\begin{annexitem}{annex:p15.stirling-poisson-normal}{thm:p15-stirling-poisson-normal}{The source Poisson calculation with uniform remainder and summation}
\begin{proof}
For \(k\geq1\), put \(u=(k-n)/n\) and
\(J(u)=(1+u)\log(1+u)-u\).
The exact factorial formula gives
\[
 \mathbb P(T_n=k)
 =\frac{\exp\{-nJ(u)-r_k\}}{\sqrt{2\pi n}\sqrt{1+u}}.
\]
If \(|k-n|\leq A\sqrt n\), then \(|u|\leq1/2\)
and \(k\geq n/2\geq1\) for all sufficiently large \(n\).
Since \(J(0)=J'(0)=0\), \(J''(0)=1\) and
\(J'''(u)=-(1+u)^{-2}\), Taylor's integral remainder gives
\[
 |J(u)-u^2/2|\leq\tfrac23|u|^3\quad(|u|\leq1/2).
\]
Also \(|\log(1+u)|\leq2|u|\) and \(0<r_k<1/(6n)\).
For \(z=(k-n)/\sqrt n\), the logarithm of the ratio to
\((2\pi n)^{-1/2}e^{-z^2/2}\) is consequently bounded
in absolute value by
\[
 \frac{2A^3/3+A}{\sqrt n}+\frac1{6n}.
\]
This proves the required uniform local approximation.
For a moving integer \(k(n)\) whose scaled argument
tends to \(x\), it also proves exactly the source
\(\sqrt{2\pi n}\)-normalized limit.
The exact \(k=0\) mass is \(e^{-n}\); no value \(r_0\)
is introduced.

The scaled lattice now has mesh \(1/\sqrt n\). The
uniform estimate on any fixed \([a,b]\) makes its total
summation error tend to zero, by the same bound of
interval length plus two mesh lengths times the uniform
error. Uniform continuity of \(\phi\) gives the Riemann
sum limit. The exact integer thresholds are obtained
by solving \(n+a\sqrt n\leq k\leq n+b\sqrt n\) and
intersecting with \(k\geq0\).
Since \(\mathbb E(T_n-n)^2=n\), Chebyshev gives
\[
 \mathbb P(|T_n-n|/\sqrt n>A)\leq A^{-2}.
\]
This proves the CDF and infinite-endpoint conclusions by
finite cutoffs. At a fixed finite scaled endpoint the
local estimate bounds its mass by a constant times
\(n^{-1/2}\), so the strict endpoint and singleton
claims follow. The source states the local limit as
Exercise~3.1.2 and indicates this interval consequence;
both calculations have now been supplied.
\end{proof}
\end{annexitem}

\begin{annexitem}{annex:p15.stirling-rademacher-tail-rate}{thm:p15-stirling-rademacher-tail-rate}{The source point-mass and geometric-tail argument}
\begin{proof}
Fix \(0<a<1\), and put \(k_n=\lceil na\rceil\).
Then \(k_n/n\to a\) and \(1\leq k_n<n\) for all
sufficiently large \(n\).
The exact Stirling factorization, now at \(u_n=k_n/n\),
gives
\[
 \frac1{2n}\log\mathbb P(S_{2n}=2k_n)
 =-\frac12H(u_n)
  -\frac{\log(\pi n)+\log(1-u_n^2)}{4n}
  +\frac{R_{n,k_n}}{2n}
 \longrightarrow-\gamma(a).
\]
Here \(u_n\) remains in a compact subinterval of \((0,1)\),
so its logarithmic and remainder terms divided by \(n\)
vanish, and \(H(a)/2=\gamma(a)\).
No central-window normal approximation is used at this
linear-deviation scale.

For \(k\geq k_n\), the successive allowed-lattice ratio is
\[
 \frac{\mathbb P(S_{2n}=2k+2)}
      {\mathbb P(S_{2n}=2k)}
    =\frac{n-k}{n+k+1}
    \leq\frac{1-a}{1+a}=:q_a<1.
\]
The event with real threshold \(2na\) starts at \(2k_n\).
Summing its masses therefore lies between the first term
and that term times the geometric factor \((1-q_a)^{-1}\).
Taking logarithms, dividing by \(2n\), and using that this
factor is fixed proves the source tail rate.

For \(a\leq0\), the event contains \(\{S_{2n}\geq0\}\),
whose probability is at least \(1/2\) by symmetry.
Its logarithm divided by \(2n\) tends to zero by squeezing
between \(-(\log2)/(2n)\) and zero.
For \(a=1\), the event requires all signs to be positive,
so its probability is \(2^{-2n}\) and its rate is
\(-\log2\). For \(a>1\) the event is empty and the
stated extended logarithmic convention gives \(-\infty\).
The endpoint value \(\gamma(1)=\log2\) agrees with the
interior formula by \(t\log t\to0\) at zero.
\end{proof}
\end{annexitem}

\begin{annexitem}{annex:p15.stirling-poisson-tail-rate}{thm:p15-stirling-poisson-tail-rate}{The Poisson tail, integer threshold and endpoint rates}
\begin{proof}
For fixed \(a>1\), put \(k_n=\lceil na\rceil\).
The exact mass and Stirling's formula imply
\[
 \frac1n\log\mathbb P(T_n=k_n)
 =-1+\frac{k_n}{n}
       -\frac{k_n}{n}\log(k_n/n)
       -\frac{\log(2\pi k_n)}{2n}-\frac{r_{k_n}}n
 \longrightarrow a-1-a\log a.
\]
For \(k\geq k_n\), its successive mass ratio is
\[
 \frac{\mathbb P(T_n=k+1)}{\mathbb P(T_n=k)}
       =\frac{n}{k+1}\leq\frac1a<1.
\]
Thus the infinite tail lies between its first mass and
that mass times \((1-1/a)^{-1}\).
The logarithm of this fixed factor divided by \(n\)
tends to zero, proving the source limit and both finite
comparison inequalities.

At \(a=1\), the tail contains \(T_n=n\), whose mass is
asymptotic to \((2\pi n)^{-1/2}\) by the factorial formula.
Its logarithm divided by \(n\) tends to zero; the tail
is at most one, so its rate is zero.
For \(0<a<1\), Chebyshev gives
\(\mathbb P(T_n<na)\leq1/(n(1-a)^2)\to0\).
The tail therefore tends to one and its logarithmic
rate is zero. For \(a\leq0\), \(T_n\geq na\) surely
because \(T_n\) is nonnegative. These cases do not use
the geometric factor at its singular endpoint.
\end{proof}
\end{annexitem}

\begin{annexitem}{annex:p15.stirling-general-binomial}{thm:p15-stirling-general-binomial}{A separately stated fixed-parameter extension}
\begin{proof}
Fix \(0<p<1\) and \(q=1-p\). For \(0<k<N\), put \(t=k/N\).
Define
\[
 D(t\Vert p)=t\log(t/p)+(1-t)\log((1-t)/q)
                                      \quad(0<t<1).
\]
The exact Stirling formula and the binomial mass give
\[
 \mathbb P(B_N=k)
 =\frac{\exp\{-ND(t\Vert p)+r_N-r_k-r_{N-k}\}}
              {\sqrt{2\pi N t(1-t)}}.
\]
It satisfies \(D(p\Vert p)=0\),
\(\partial_tD(p\Vert p)=0\), and
\(\partial_t^2D(t\Vert p)=1/(t(1-t))\).
Take \(\delta=\min(p,q)/2\) and the interval
\([p-\delta,p+\delta]\subset(0,1)\).
There the third derivative
\(-t^{-2}+(1-t)^{-2}\) has absolute value at most
\(M_p=4/p^2+4/q^2\).
Taylor's integral remainder therefore gives
\[
 \left|D(p+h\Vert p)-\frac{h^2}{2pq}\right|
                           \leq M_p|h|^3/6
\]
for every \(p+h\) in that interval.

For \(|k-Np|\leq A\sqrt{Npq}\), one has
\(|h|=|t-p|\leq A\sqrt{pq/N}\).
All such \(t\) lie in the fixed interior interval
eventually. It follows uniformly that
\[
 ND(t\Vert p)=\frac{(k-Np)^2}{2Npq}+o(1),\qquad
 r_N-r_k-r_{N-k}=o(1),\qquad
 \frac{t(1-t)}{pq}\to1.
\]
The Taylor remainder is explicitly at most
\(M_pA^3(pq)^{3/2}/(6\sqrt N)\);
the factorial remainder is bounded by a constant
depending on \(p\) times \(1/N\).
Substitution proves the stated uniform local limit.
In particular no approximation is assumed uniform in
\(p\) approaching zero or one.

The scaled lattice has mesh \(1/\sigma_N\to0\).
Uniform continuity of \(\phi\) and the preceding uniform
local bound give the interval limits by Riemann sums,
with exact endpoints obtained by ceiling and floor as
stated. The variance identity yields the uniform tail
bound
\(\mathbb P(|B_N-Np|/\sigma_N>A)\leq A^{-2}\).
Finite cutoffs therefore prove convergence of every CDF
value to the continuous \(\Phi\), which is convergence
in distribution. This proof is a fixed-parameter
extension of the factorial method and does not import
an existing normal-limit result.
At \(p=0\) all trials fail and at \(p=1\) all succeed.
Then \(B_N=Np\) surely and the variance and stated
normalizer are zero, proving the separate degeneracy
qualification.
\end{proof}
\end{annexitem}

\begin{annexitem}{annex:p15.stirling-worked-boundaries}{thm:p15-stirling-worked-boundaries}{Computed examples and precise failures}
\begin{proof}
By the exact factorial formula, \(A_n/n!=e^{-r_n}\).
Since \(r_n>0\) strictly decreases to zero, the relative
error \(1-e^{-r_n}\) is positive and strictly decreases
to zero. Nevertheless \(e^{r_n}\geq1+r_n\), so
\[
 n!-A_n=n!(1-e^{-r_n})
       \geq n!\frac{r_n}{1+r_n}
       >\frac{n!}{12n+2}.
\]
For \(n\geq4\), at least \(\lfloor n/2\rfloor\) factors
in \(n!\) are at least \(n/2\). The resulting lower
bound grows faster than \(12n+2\), proving divergence
of the absolute error. This gives the source distinction
between absolute and percentage error.

The exact expressions at \(n=5\) are
\[
 A_5=\sqrt{2\pi}\,5^{11/2}e^{-5}
       =118.0191679575\ldots,\qquad
 A_5e^{1/60}=120.0026370861\ldots.
\]
The latter is a strict upper bound for \(120\) by
Robbins's inequality. Direct substitution in the two
rate functions gives
\[
 -\gamma(1/2)
 =-\tfrac12\{\tfrac32\log(3/2)+\tfrac12\log(1/2)\}
 =-0.1308120359\ldots,
\]
and \(1-2\log2=-0.3862943611\ldots\).

If all signs equal one common fair sign \(\varepsilon\),
then \(S_N/\sqrt N=\sqrt N\,\varepsilon\).
For \(N>1\) it has no probability in \([-1,1]\),
whereas \(\int_{-1}^1\phi>0\).
This explicitly fails the normal limit under dependence.
For independent signs, the all-positive endpoint has
probability \(4^{-n}\) at time \(2n\), giving exactly
the displayed ratio \(\sqrt{\pi n}(e/4)^n\).
It tends to zero because \(e<3<4\): the exponential
series gives \(e=2+\sum_{j\geq2}1/j!<3\) by comparison
with \(\sum_{j\geq2}2^{-(j-1)}=1\).
An exponential with ratio below one dominates the
factor \(\sqrt n\), for example by its consecutive
ratio tending to \(e/4<1\).
Thus the relative central approximation cannot be
extended uniformly to the entire support.

For \(p_N=1/N\), \(N\geq2\), the binomial mean is one
and variance \(1-1/N\). At \(B_N=0\) the standardized
value is \(-1/\sqrt{1-1/N}\to-1\), with probability
\((1-1/N)^N\to e^{-1}\) by the proved scalar product
limit. Let \(K=[-1-1/10,-1+1/10]\).
For all sufficiently large \(N\), that atom lies in \(K\),
so the lower limit of the probability of \(K\) is at
least \(e^{-1}\).
The standard normal probability of \(K\) is at most
\((1/5)/\sqrt{2\pi}<1/5<e^{-1}\), since \(e<3\).
If these laws converged to the standard normal law,
the closed-set portmanteau bound would contradict
this inequality. Thus arbitrary changing probabilities
are outside the stated fixed-\(p\) result.
\end{proof}
\end{annexitem}
